% 第 XXI 讲尾部；冻结底本源行 3045--5057；由总编与前部合并。
\section{根数据的态射}\label{sga3:XXI:6}
\renewcommand{\thetheorem}{\thesubsection.\arabic{theorem}}

\subsection{定义}\label{sga3:XXI:6.1}

设
\[
 \cR=(M,M^\ast,R,R^\ast),
 \qquad
 \cR'=(M',M'^\ast,R',R'^\ast)
\]
为两个根数据。设 \(f:M'\to M\) 为线性映射，\({}^tf:M^\ast\to M'^\ast\) 为其转置。

\begin{definition}\label{sga3:XXI:6.1.1}
若 \(f\) 诱导 \(R'\) 到 \(R\) 的双射，且 \({}^tf\) 诱导 \(R^\ast\) 到 \(R'^\ast\) 的双射，则称映射 \(f\) 为 \(\cR'\) 到 \(\cR\) 的一个\emph{态射}，记为
\[
 f:\cR'\longrightarrow\cR.
\]
\end{definition}

此时 \({}^tf\) 为对偶根数据之间的态射：
\[
 {}^tf:\cR^\ast\longrightarrow\cR'^\ast.
\]
若 \(f:\cR'\to\cR\) 为态射，且对 \(\alpha'\in R'\) 有 \(\alpha=f(\alpha')\)，则
\[
 \alpha'^\ast={}^tf(\alpha^\ast).
\]
事实上，若 \(p\) 与 \(p'\) 分别表示 \(\cR\) 与 \(\cR'\) 的 \ref{sga3:XXI:1.2.1} 中的映射，则
\[
 p'={}^tf\circ p\circ f,
\]
由此立即得到断言。读者可自行验证下述各断言，它们几乎都是直接的。

\begin{proposition}\label{sga3:XXI:6.1.2}
设 \(f:\cR'\to\cR\) 为根数据的态射。若 \(\alpha'\in R'\)，且 \(\alpha=f(\alpha')\)，则
\[
 \alpha'^\ast={}^tf(\alpha^\ast).
\]
此外，\(f\) 诱导同构
\[
 R'\xrightarrow{\sim}R,\qquad
 \Gamma_0(R')\xrightarrow{\sim}\Gamma_0(R),\qquad
 \VR(R')\xrightarrow{\sim}\VR(R),
\]
而 \({}^tf\) 诱导同构
\[
 R^\ast\xrightarrow{\sim}R'^\ast,\qquad
 \Gamma_0(R^\ast)\xrightarrow{\sim}\Gamma_0(R'^\ast),\qquad
 \VR(R^\ast)\xrightarrow{\sim}\VR(R'^\ast).
\]
最后一个同构是由 \(f\) 诱导的相应态射的转置。映射
\[
 s_{\alpha'}\longmapsto s_{f(\alpha')}
\]
延拓为一个同构
\[
 W(\cR')\xrightarrow{\sim}W(\cR),
\]
它与这两个群在 \ref{sga3:XXI:1.1.13} 中各集合上的作用相容。
\end{proposition}

\begin{proposition}\label{sga3:XXI:6.1.3}
映射
\[
 \Delta'\longmapsto f(\Delta'),\qquad
 R'_+\longmapsto f(R'_+),\qquad
 C'\longmapsto({}^tf\otimes\mathbb R)^{-1}(C')
\]
定义了 \(\cR'\) 与 \(\cR\) 的单根系、正根系及 Weyl 腔之间的双射对应。这些对应与 Weyl 群的作用以及对应
\[
 \mathscr{S}(R_+)\longleftrightarrow R_+
 \longleftrightarrow\mathscr{C}(R_+)
\]
相容。
\end{proposition}

\begin{lemma}\label{sga3:XXI:6.1.4}
态射可复合。态射 \(f:\cR'\to\cR\) 为同构，当且仅当线性映射 \(f:M'\to M\) 为双射。
\end{lemma}

\subsection{同源}\label{sga3:XXI:6.2}

\begin{definition}\label{sga3:XXI:6.2.1}
若根数据的态射 \(f:\cR'\to\cR\) 满足 \(f:M'\to M\) 为单射且余核有限，则称它为一个\emph{同源}。
\end{definition}

若 \(f\) 为同源，则 \({}^tf\) 为同源。

\begin{definition}\label{sga3:XXI:6.2.2}
设 \(f:\cR'\to\cR\) 为同源。令
\[
 K(f)=\operatorname{Coker}\bigl(M'\xrightarrow{f}M\bigr).
\]
\end{definition}

\begin{lemma}\label{sga3:XXI:6.2.3}
存在自然配对
\[
 K(f)\times K({}^tf)\longrightarrow\mathbb Q/\mathbb Z,
\]
使这两个有限群互为对偶。
\end{lemma}

\begin{proof}
这是经典结果。
\end{proof}

\begin{lemma}\label{sga3:XXI:6.2.4}
若 \(f:\cR'\to\cR\) 为态射，则
\[
 \rgss(\cR')=\rgss(\cR).
\]
若进一步 \(f\) 为同源，则
\[
 \rgred(\cR')=\rgred(\cR).
\]
\end{lemma}

\begin{proof}
这是直接的。
\end{proof}

\begin{lemma}\label{sga3:XXI:6.2.5}
半单根数据之间的每个态射均为同源。
\end{lemma}

\begin{proof}
映射 \(f\) 必须诱导 \(V'=\VR(R')\) 到 \(V=\VR(R)\) 的同构。
\end{proof}

若 \(\cR'\) 与 \(\cR\) 半单，则每个同源 \(f:\cR'\to\cR\) 定义交换图：
\begin{center}
\begin{tikzcd}[column sep=huge,row sep=large]
\Gamma_0(R') \arrow[r,"\sim"] \arrow[d]
  & \Gamma_0(R) \arrow[d] \\
M' \arrow[r,"f"'] & M
\end{tikzcd}
\end{center}
若 \(M=\Gamma_0(R)\)，则 \(f\) 必为同构。

\begin{definition}\label{sga3:XXI:6.2.6}
若根数据满足
\[
 M=\Gamma_0(R),
 \qquad\text{分别为}\qquad
 M^\ast=\Gamma_0(R^\ast),
\]
则称它为\emph{伴随的}，分别为\emph{单连通的}。
\end{definition}

因此伴随或单连通根数据均半单。此外，\(\cR\) 伴随，分别为单连通，当且仅当 \(\cR^\ast\) 单连通，分别为伴随。上述结果给出下述命题。

\begin{proposition}\label{sga3:XXI:6.2.7}
设 \(\cR\) 为半单根数据。下述条件等价：
\begin{enumerate}
\item[(i)] \(\cR\) 伴随，分别为单连通；
\item[(ii)] 每个同源 \(\cR'\to\cR\)，分别为 \(\cR\to\cR'\)，都是同构。
\end{enumerate}
\end{proposition}

\begin{proposition}\label{sga3:XXI:6.2.8}
设 \(\cR\) 为伴随，分别为单连通的根数据。每个不可分根，分别为余根，都是 \(M\)，分别为 \(M^\ast\)，中的不可分元素。
\end{proposition}

\begin{proof}
由 \ref{sga3:XXI:3.3.5}，每个不可分根属于 \(\Gamma_0(R)\) 的某个基。
\end{proof}

\subsection{根基与余根基}\label{sga3:XXI:6.3}

设 \(\cR\) 为根数据。令
\[
 N=\{x\in M\mid(\alpha^\ast,x)=0
       \text{ 对每个 }\alpha^\ast\in R^\ast\},
\qquad
 N^\ast=M^\ast/\bigl(\VR(R^\ast)\cap M^\ast\bigr).
\]

\begin{lemma}\label{sga3:XXI:6.3.1}
考虑典范态射
\[
 N\longrightarrow M,
 \qquad
 M^\ast\longrightarrow N^\ast.
\]
它们互为转置，且 \(N^\ast\) 与 \(N\) 的对偶相认同。
\end{lemma}

\begin{proof}
这直接由 \ref{sga3:XXI:1.2.5} 得到。
\end{proof}

\begin{definition}\label{sga3:XXI:6.3.2}
\(\cR\) 的\emph{余根基}，记为 \(\operatorname{corad}(\cR)\)，是平凡根数据
\[
 \operatorname{corad}(\cR)=(N,N^\ast,\varnothing,\varnothing).
\]
\end{definition}

若
\[
 \cR^0=(M,M^\ast,\varnothing,\varnothing),
\]
即一个平凡根数据，则有态射
\[
 \operatorname{corad}(\cR)\longrightarrow\cR^0.
\]

\begin{definition}\label{sga3:XXI:6.3.3}
\(\cR\) 的\emph{根基}，记为 \(\operatorname{rad}(\cR)\)，是平凡根数据
\footnote{编者注：等价地，若
\[
 P=\{y\in M^\ast\mid(y,\alpha)=0\text{ 对每个 }\alpha\in R\},
 \qquad
 P^\ast=M/\bigl(\VR(R)\cap M\bigr),
\]
则 \(\operatorname{rad}(\cR)=(P^\ast,P,\varnothing,\varnothing)\)。}
\[
 \operatorname{rad}(\cR)=\operatorname{corad}(\cR^\ast)^\ast.
\]
\end{definition}

因此有图
\begin{center}
\begin{tikzcd}[column sep=huge,row sep=large]
\operatorname{corad}(\cR) \arrow[rr] \arrow[dr]
  && \operatorname{rad}(\cR) \\
& \cR^0 \arrow[ur] &
\end{tikzcd}
\end{center}
其转置为 \(\cR^\ast\) 对应的图。

\begin{lemma}\label{sga3:XXI:6.3.4}
典范态射
\[
 u:\operatorname{corad}(\cR)\longrightarrow\operatorname{rad}(\cR)
\]
为同源。
\end{lemma}

\begin{definition}\label{sga3:XXI:6.3.5}
令
\[
 N(\cR)=K(u)
 =
 M\Big/
 \left(
   \bigl(\VR(R)\cap M\bigr)
   +\bigcap_{\alpha\in R}\ker(\alpha^\ast)
 \right).
\]
则有典范配对
\[
 N(\cR)\times N(\cR^\ast)\longrightarrow\mathbb Q/\mathbb Z.
\]
\end{definition}

\begin{lemma}\label{sga3:XXI:6.3.6}
有
\[
 \rgred\bigl(\operatorname{rad}(\cR)\bigr)
 =
 \rgred\bigl(\operatorname{corad}(\cR)\bigr)
 =
 \rgred(\cR)-\rgss(\cR),
\]
且下述条件等价：
\begin{enumerate}
\item[(i)] \(\cR\) 半单；
\item[(ii)] \(\operatorname{rad}(\cR)=0\)；
\item[(iii)] \(\operatorname{corad}(\cR)=0\)。
\end{enumerate}
\end{lemma}
% 底本源行 3045--3354 译毕。

\subsection{根数据的积}\label{sga3:XXI:6.4}

\begin{definition}\label{sga3:XXI:6.4.1}
设
\[
 \cR=(M,M^\ast,R,R^\ast),
 \qquad
 \cR'=(M',M'^\ast,R',R'^\ast)
\]
为两个根数据。\(\cR\) 与 \(\cR'\) 的\emph{积根数据}，记为
\[
 \cR''=\cR\times\cR',
\]
是根数据 \((M'',M''^\ast,R'',R''^\ast)\)，其中
\[
\begin{aligned}
 M''&=M\times M',
 &M''^\ast&=M^\ast\times M'^\ast,\\
 R''&=(R\times0)\cup(0\times R'),
 &R''^\ast&=(R^\ast\times0)\cup(0\times R'^\ast),
\end{aligned}
\]
而映射 \(\alpha\mapsto\alpha^\ast\) 为显然的映射。
\end{definition}

\begin{proposition}\label{sga3:XXI:6.4.2}
在上述假设下，有典范同构
\[
 \Gamma_0(R'')\simeq\Gamma_0(R)\times\Gamma_0(R'),
 \qquad
 \VR(R'')\simeq\VR(R)\times\VR(R'),
\]
\[
 W(\cR'')\simeq W(\cR)\times W(\cR'),
\]
其余附属构造也有类似结果。此外，
\[
 \rgred(\cR'')=\rgred(\cR)+\rgred(\cR'),
 \qquad
 \rgss(\cR'')=\rgss(\cR)+\rgss(\cR').
\]
还有根数据的典范同构
\[
 (\cR\times\cR')^\ast\simeq\cR^\ast\times\cR'^\ast.
\]
\end{proposition}

上述定义立即推广到任意有限个因子的积。

\begin{proposition}\label{sga3:XXI:6.4.3}
设
\[
 \cR=\cR_1\times\cR_2\times\cdots\times\cR_n
\]
为根数据的积。下述条件等价：
\begin{enumerate}
\item[(i)] \(\cR\) 半单，分别为单连通、伴随、约化；
\item[(ii)] 每个 \(\cR_i\) 均半单，分别为单连通、伴随、约化。
\end{enumerate}
\end{proposition}

考虑下述特殊情形。设 \(\cR_0\) 为平凡根数据，\(\cR_1\) 为半单根数据。此时有交换图
\begin{center}
\begin{tikzcd}[column sep=large,row sep=large]
& \cR_0\times\cR_1 \arrow[dr] & \\
\cR_1 \arrow[ur] \arrow[rr,"\operatorname{id}"]
  && \cR_1
\end{tikzcd}
\end{center}

\begin{lemma}\label{sga3:XXI:6.4.4}
有典范同构
\begin{center}
\begin{tikzcd}[column sep=huge,row sep=large]
\operatorname{corad}(\cR_0\times\cR_1)
  \arrow[rr,"\sim"] \arrow[dr,"\sim"']
  && \operatorname{rad}(\cR_0\times\cR_1) \\
& \cR_0 \arrow[ur,"\sim"'] &
\end{tikzcd}
\end{center}
特别地，
\[
 N(\cR_0\times\cR_1)=0.
\]
\end{lemma}

后面将看到，反过来，若 \(N(\cR)=0\)，则 \(\cR\) 为一个半单根数据与一个平凡根数据的积。

\subsection{诱导根数据与余诱导根数据}\label{sga3:XXI:6.5}

设 \(\cR=(M,M^\ast,R,R^\ast)\) 为根数据，\(N\subset M\) 为包含各根的子群，即
\[
 \Gamma_0(R)\subset N\subset M.
\]
典范线性映射 \(i_N:N\to M\) 经转置给出线性映射
\[
 {}^ti_N:M^\ast\longrightarrow N^\ast.
\]
令
\[
 R_N=R,
 \qquad
 R_N^\ast={}^ti_N(R^\ast).
\]

\begin{lemma}\label{sga3:XXI:6.5.1}
四元组
\[
 \cR_N=(N,N^\ast,R_N,R_N^\ast)
\]
是根数据，且 \(i_N\) 为态射。
\end{lemma}

\begin{proof}
首先，\({}^ti_N\) 诱导 \(R^\ast\) 到 \(R_N^\ast\) 的同构。事实上，若 \(\alpha,\beta\in R\)，且
\[
 {}^ti_N(\alpha^\ast)={}^ti_N(\beta^\ast),
\]
则对每个 \(x\in N\)，特别是对 \(x\in R\)，有
\[
 (\alpha^\ast,x)=(\beta^\ast,x).
\]
因此由 \ref{sga3:XXI:1.1.4}，\(\alpha=\beta\)。其余部分不难得到。
\end{proof}

\begin{definition}\label{sga3:XXI:6.5.2}
称根数据 \(\cR_N\) 为\emph{由 \(\cR\) 在 \(N\) 上诱导的根数据}。
\end{definition}

\begin{lemma}\label{sga3:XXI:6.5.3}
设 \(f:\cR'\to\cR\) 为态射，令
\[
 N=f(M')\subset M.
\]
则 \(f\) 唯一地经过 \(i_N\) 分解。
\end{lemma}

特别地，在同构意义下，同源 \(\cR'\to\cR\) 与
\[
 M/\Gamma_0(R)
\]
的有限指数子群一一对应。这加强了 \ref{sga3:XXI:6.2.7}。

现设 \(N^\ast\subset M^\ast\) 为包含 \(R^\ast\) 的子群。\emph{由 \(\cR\) 在 \(N^\ast\) 上余诱导的根数据}定义为
\[
 \cR^{N^\ast}=(\cR^\ast_{N^\ast})^\ast.
\]
有典范态射
\[
 p^{N^\ast}:\cR\longrightarrow\cR^{N^\ast}.
\]

\begin{lemma}\label{sga3:XXI:6.5.4}
设 \(f:\cR'\to\cR\) 为态射。存在子群 \(N\subset M\) 与 \(N'^\ast\subset M'^\ast\)，使 \(f\) 分解为
\begin{center}
\begin{tikzcd}[column sep=huge,row sep=large]
\cR' \arrow[r,"f"] \arrow[d,"p^{N'^\ast}"']
  & \cR \\
\cR'^{\,N'^\ast} \arrow[r,"f_0","\sim"']
  & \cR_N \arrow[u,"i_N"']
\end{tikzcd}
\end{center}
其中 \(f_0\) 为同构。
\end{lemma}

\begin{proof}
与 \ref{sga3:XXI:6.5.3} 一样，取 \(N=f(M')\)。所得态射 \(M'\to N\) 为满射，故其转置为单射。取此转置的像为 \(N'^\ast\)。
\end{proof}

现处理几个特殊情形。若 \(N=\Gamma_0(R)\)，记
\[
 \cR_N=\operatorname{ad}(\cR).
\]
若 \(N=\VR(R)\cap M\)，记
\[
 \cR_N=\operatorname{ss}(\cR).
\]
因此有图
\[
 \operatorname{ad}(\cR)\longrightarrow
 \operatorname{ss}(\cR)\longrightarrow\cR.
\]
令
\[
 \operatorname{der}(\cR)=\operatorname{ss}(\cR^\ast)^\ast,
 \qquad
 \operatorname{sc}(\cR)=\operatorname{ad}(\cR^\ast)^\ast.
\]
由对偶性，有图
\[
 \cR\longrightarrow
 \operatorname{der}(\cR)\longrightarrow
 \operatorname{sc}(\cR).
\]

\begin{proposition}\label{sga3:XXI:6.5.5}
\begin{enumerate}
\item[(i)] 图
\begin{center}
\begin{tikzcd}[column sep=large,row sep=large]
\operatorname{ad}(\cR) \arrow[r]
  & \operatorname{ss}(\cR) \arrow[rr] \arrow[dr]
  && \operatorname{der}(\cR) \arrow[r]
  & \operatorname{sc}(\cR) \\
&& \cR \arrow[ur] &&
\end{tikzcd}
\end{center}
第一行中的四个根数据均半单，三个态射均为同源。
\item[(ii)] 数据 \(\operatorname{ad}(\cR)\) 伴随；且 \(\cR\) 伴随，当且仅当 \(\operatorname{ad}(\cR)\to\cR\) 为同构。
\item[(iii)] 数据 \(\operatorname{sc}(\cR)\) 单连通；且 \(\cR\) 单连通，当且仅当 \(\cR\to\operatorname{sc}(\cR)\) 为同构。
\item[(iv)] 下述条件等价：
\begin{enumerate}
\item[(a)] \(\cR\) 半单；
\item[(b)] \(\operatorname{ss}(\cR)\to\cR\) 为同构；
\item[(c)] \(\cR\to\operatorname{der}(\cR)\) 为同构。
\end{enumerate}
\end{enumerate}
\end{proposition}

考虑态射
\[
 \operatorname{ss}(\cR)\longrightarrow\operatorname{der}(\cR).
\]
由定义立即得到下述结果。

\begin{lemma}\label{sga3:XXI:6.5.6}
设
\[
 h:\operatorname{ss}(\cR)\longrightarrow\operatorname{der}(\cR)
\]
为典范同源。则
\[
 K(h)\simeq N(\cR).
\]
\end{lemma}

\SGARefTarget{sga3:XXI:6.5.7}
考虑诱导根数据的其他特殊情形。令
\[
 N=
 \{x\in M\mid(\alpha^\ast,x)=0\text{ 对每个 }\alpha\in R\}
 \times\Gamma_0(R).
\]
由 \ref{sga3:XXI:1.2.5}，这个和是直和。于是 \(\cR_N\) 可与
\[
 \operatorname{ad}(\cR)\times\operatorname{corad}(\cR)
\]
相认同。可将 \(\Gamma_0(R)\) 换为 \(\VR(R)\cap M\)，再将这两个构造对偶化。这给出下述根数据图：
\begin{center}
\resizebox{\textwidth}{!}{%
\begin{tikzpicture}[
  >=stealth,
  every node/.style={font=\small},
  x=1cm,y=1cm
]
\node (at) at (0,2.1) {$\operatorname{ad}(\cR)$};
\node (st) at (4,2.1) {$\operatorname{ss}(\cR)$};
\node (dt) at (10,2.1) {$\operatorname{der}(\cR)$};
\node (ct) at (14,2.1) {$\operatorname{sc}(\cR)$};

\node (am) at (0,0) {$\operatorname{ad}(\cR)\times\operatorname{corad}(\cR)$};
\node (sm) at (4,0) {$\operatorname{ss}(\cR)\times\operatorname{corad}(\cR)$};
\node (rm) at (7,0) {$\cR$};
\node (dm) at (10,0) {$\operatorname{der}(\cR)\times\operatorname{rad}(\cR)$};
\node (cm) at (14,0) {$\operatorname{sc}(\cR)\times\operatorname{rad}(\cR)$};

\node (ab) at (0,-2.1) {$\operatorname{ad}(\cR)$};
\node (sb) at (4,-2.1) {$\operatorname{ss}(\cR)$};
\node (db) at (10,-2.1) {$\operatorname{der}(\cR)$};
\node (cb) at (14,-2.1) {$\operatorname{sc}(\cR)$};

\draw[->] (at) -- (st);
\draw[->] (st) -- (dt);
\draw[->] (dt) -- (ct);
\draw[->] (am) -- (sm);
\draw[->] (sm) -- (rm);
\draw[->] (rm) -- (dm);
\draw[->] (dm) -- (cm);
\draw[->] (ab) -- (sb);
\draw[->] (sb) -- (db);
\draw[->] (db) -- (cb);

\draw[->] (at) -- (am);
\draw[->] (am) -- (ab);
\draw[->] (st) -- (sm);
\draw[->] (sm) -- (sb);
\draw[->] (dt) -- (dm);
\draw[->] (dm) -- (db);
\draw[->] (ct) -- (cm);
\draw[->] (cm) -- (cb);
\end{tikzpicture}
}
\end{center}
立即可验证此图交换，并且在显然的意义下自对偶。各水平态射均为同源，各竖直箭头的复合均为恒等。

\begin{lemma}\label{sga3:XXI:6.5.8}
设 \(h_1\) 与 \(h_2\) 为典范同源
\[
 \operatorname{ss}(\cR)\times\operatorname{corad}(\cR)
 \xrightarrow{h_1}\cR
 \xrightarrow{h_2}
 \operatorname{der}(\cR)\times\operatorname{rad}(\cR).
\]
则
\[
 K(h_1)\simeq K(h_2)\simeq N(\cR).
\]
\end{lemma}

\begin{proof}
这直接由定义得到。
\end{proof}

\begin{corollary}\label{sga3:XXI:6.5.9}
对根数据 \(\cR\)，下述条件等价：
\begin{enumerate}
\item[(i)] \(N(\cR)=0\)，即 \(\operatorname{corad}(\cR)\to\operatorname{rad}(\cR)\) 为同构；
\item[(ii)] 典范映射 \(\operatorname{ss}(\cR)\to\operatorname{der}(\cR)\) 为同构；
\item[(iii)] \(h_1\) 为同构；
\item[(iv)] \(h_2\) 为同构；
\item[(v)] \(\cR\) 为一个半单根数据与一个平凡根数据的积。
\end{enumerate}
\end{corollary}

\begin{corollary}\label{sga3:XXI:6.5.10}
对每个根数据 \(\cR\)，存在同源
\[
 \operatorname{ad}(\cR)\times\cR_0
 \longrightarrow\cR\longrightarrow
 \operatorname{sc}(\cR)\times\cR_0,
\]
其中 \(\cR_0\) 为秩是 \(\rgred(\cR)-\rgss(\cR)\) 的平凡根数据。
\end{corollary}

\begin{lemma}\label{sga3:XXI:6.5.11}
设 \(\cR\) 为根数据，\(\Delta\) 为单根系，\(\Delta'\subset\Delta\)。考虑 \ref{sga3:XXI:3.4.7} 中的根数据
\[
 \cR_{\Delta'}=(M,M^\ast,R_{\Delta'},R^\ast_{\Delta'}).
\]
\begin{enumerate}
\item[(i)] 若 \(\cR\) 单连通，则 \(\operatorname{der}(\cR_{\Delta'})\) 单连通。
\item[(ii)] 若 \(\cR\) 伴随，则 \(\operatorname{ss}(\cR_{\Delta'})\) 伴随。
\end{enumerate}
\end{lemma}

\begin{proof}
由对偶性，两项断言等价。第二项化为
\[
 M\cap\VR(R_{\Delta'})=\Gamma_0(R_{\Delta'}).
\]
当 \(M=\Gamma_0(R)\) 时，两边均为 \(M\) 中由 \(\Delta'\) 生成的子群。
\end{proof}
% 底本源行 3355--3758 译毕。

\subsection{权}\label{sga3:XXI:6.6}

\begin{definition}\label{sga3:XXI:6.6.1}
对根数据 \(\cR\)，令
\footnote{编者注：为避免与 \(\Gamma_0(R)\) 混淆，记号 \(\Gamma(\cR)\) 已改为 \(\Lambda(\cR)\)。}
\[
 \Lambda(\cR)=
 \{x\in\VR(R)\mid(\alpha^\ast,x)\in\mathbb Z
   \text{ 对每个 }\alpha^\ast\in R^\ast\}.
\]
\(\Lambda(\cR)\) 的元素称为 \(\cR\) 的\emph{权}。\(\cR^\ast\) 的权称为 \(\cR\) 的\emph{余权}。
\end{definition}

有
\[
 \Gamma_0(R)\subset\Lambda(\cR),
\]
且 \(\Lambda(\cR)\) 在 \(W(\cR)\) 下稳定。

\begin{lemma}\label{sga3:XXI:6.6.2}
双线性映射
\[
 V^\ast\times V\longrightarrow\mathbb Q
\]
诱导对偶配对
\[
 \Gamma_0(R^\ast)\times\Lambda(\cR)\longrightarrow\mathbb Z.
\]
\end{lemma}

\begin{proof}
直接可得。
\end{proof}

\begin{corollary}\label{sga3:XXI:6.6.3}
设
\[
 \Delta^\ast=(\alpha_1^\ast,\ldots,\alpha_n^\ast)
\]
为单余根系，且由
\[
 (\alpha_i^\ast,p_j)=\delta_{ij}
\]
定义 \(p_i\in\VR(R)\)，\(i=1,\ldots,n\)。于是
\[
 s_{\alpha_i}(p_i)=p_i-\alpha_i,
 \qquad
 s_{\alpha_i}(p_j)=p_j\quad(i\ne j).
\footnote{注意：当数据不约化时，各 \(\alpha_i\) 不一定构成单根系。}
\]
则 \(\Lambda(\cR)\) 为各 \(p_i\) 生成的自由 Abel 群。
\end{corollary}

各 \(p_i\) 称为与单余根系 \(\Delta^\ast\) 对应的\emph{基本权}。

\begin{corollary}\label{sga3:XXI:6.6.4}
对每个 \(\alpha^\ast\in\Delta^\ast\)，有
\[
 \left(\alpha^\ast,\sum_i p_i\right)=1.
\]
因此，按 \ref{sga3:XXI:3.5.1} 的记号，
\[
 \sum_i p_i=\rho_{R_+},
 \qquad
 R_+=\mathscr P\bigl(\operatorname{ind}(\Delta)\bigr).
\]
\end{corollary}

\begin{corollary}\label{sga3:XXI:6.6.5}
对每个 \(x\in\VR(R)\)，有
\[
 x=\sum_i(\alpha_i^\ast,x)p_i.
\]
\end{corollary}

由于
\[
 R^\ast\subset\Gamma_0(R^\ast),
 \qquad
 R\subset\Lambda(\cR),
\]
四元组
\[
 \bigl(\Lambda(\cR),\Gamma_0(R^\ast),R,R^\ast\bigr)
\]
为根数据。

\begin{corollary}\label{sga3:XXI:6.6.6}
典范态射
\[
 \Gamma_0(R^\ast)\longrightarrow M^\ast
\]
为
\[
 x\longmapsto\sum_i(\alpha_i^\ast,x)p_i
\]
的转置。此映射定义根数据的态射，且有交换图
\begin{center}
\begin{tikzcd}[column sep=huge,row sep=large]
\cR \arrow[rr] \arrow[dr]
  && \operatorname{sc}(\cR) \arrow[dl,"\wr"'] \\
& \bigl(\Lambda(\cR),\Gamma_0(R^\ast),R,R^\ast\bigr) &
\end{tikzcd}
\end{center}
\end{corollary}

因此，我们用 \(\cR\) 的权给出了 \(\operatorname{sc}(\cR)\) 的明确描述。同样可得到交换图
\begin{center}
\begin{tikzcd}[column sep=huge,row sep=large]
\operatorname{ad}(\cR) \arrow[rr] \arrow[d,"\wr"]
  && \cR \\
\bigl(\Gamma_0(R),\Lambda(\cR^\ast),R,R^\ast\bigr)
  \arrow[urr] &&
\end{tikzcd}
\end{center}

\begin{corollary}\label{sga3:XXI:6.6.7}
数据 \(\cR\) 单连通，当且仅当
\[
 M=\Lambda(\cR).
\]
\end{corollary}

\begin{remark}\label{sga3:XXI:6.6.8}
有
\[
 \Lambda(\cR)\cap M=\VR(R)\cap M.
\]
因此，\(\cR\) 半单，当且仅当
\[
 M\subset\Lambda(\cR).
\]
\end{remark}

第 \ref{sga3:XXI:6.5} 节的结果还给出：

\begin{corollary}\label{sga3:XXI:6.6.9}
数据 \(\cR\) 为一个单连通数据与一个平凡数据的积，当且仅当
\[
 M\supset\Lambda(\cR).
\]
\end{corollary}

现考虑典范同源
\[
 f:\operatorname{ad}(\cR)\longrightarrow\operatorname{sc}(\cR),
\]
并令 \(Z(\cR)=K(f)\)。有
\[
 Z(\cR)\simeq Z(\operatorname{sc}\cR)
 \simeq Z(\operatorname{ad}\cR).
\]

\begin{corollary}\label{sga3:XXI:6.6.10}
有典范同构
\[
 Z(\cR)=\Lambda(\cR)/\Gamma_0(R).
\]
更确切地，有 \(W(\cR)\)-模的正合列
\[
 0\longrightarrow\Gamma_0(R)\longrightarrow\Lambda(\cR)
 \longrightarrow Z(\cR)\longrightarrow0.
\]
\end{corollary}

\begin{corollary}\label{sga3:XXI:6.6.11}
有典范配对
\[
 Z(\cR^\ast)\times Z(\cR)\longrightarrow\mathbb Q/\mathbb Z,
\]
使这两个群互为对偶。
\end{corollary}

\begin{remark}\label{sga3:XXI:6.6.12}
有
\[
 Z(\cR\times\cR')\simeq Z(\cR)\times Z(\cR').
\]
特别地，考虑单连通数据 \(\cR_i\)，\(i=1,2,\ldots,n\)，以及平凡数据 \(\cR_0\)。令
\[
 \cR=\cR_0\times\cR_1\times\cdots\times\cR_n.
\]
若
\[
 \cR=(M,M^\ast,R,R^\ast),
 \qquad
 \cR_0=(M_0,M_0^\ast,\varnothing,\varnothing),
\]
则
\[
 M/\Gamma_0(R)\simeq
 M_0\times Z(\cR_1)\times\cdots\times Z(\cR_n).
\]
\end{remark}

\subsection{自同构}\label{sga3:XXI:6.7}

由 \ref{sga3:XXI:6.1.4}，\(\cR\) 的一个自同构是 \(M\) 的一个自同构 \(u\)，使
\[
 u(R)=R,
 \qquad
 {}^tu(R^\ast)=R^\ast.
\]
特别地，每个 \(w\in W(\cR)\) 定义 \(\cR\) 的一个自同构。

\begin{lemma}\label{sga3:XXI:6.7.1}
群 \(W(\cR)\) 为 \(\operatorname{Aut}(\cR)\) 的正规子群。更确切地，若 \(u\in\operatorname{Aut}(\cR)\) 且 \(\alpha\in R\)，则
\[
 us_\alpha u^{-1}=s_{u(\alpha)}.
\]
\end{lemma}

\begin{proof}
证明与 \ref{sga3:XXI:1.2.10} 相同。
\end{proof}

\begin{proposition}\label{sga3:XXI:6.7.2}
设 \(\Delta\) 为单根系，令
\[
 E_\Delta(\cR)=
 \{u\in\operatorname{Aut}(\cR)\mid u(\Delta)=\Delta\}.
\]
则 \(\operatorname{Aut}(\cR)\) 为 \(W(\cR)\) 与 \(E_\Delta(\cR)\) 的半直积，其中前者为正规子群。
\end{proposition}

\begin{proof}
这立即由下述事实得到：\(W(\cR)\) 在各单根系上单可迁地作用；且若 \(\Delta\) 为 \(\cR\) 的单根系，则对 \(\cR\) 的每个自同构 \(u\)，\(u(\Delta)\) 都是单根系。
\end{proof}

后面将在约化不可约根数据的情形下给出 \(E_\Delta(\cR)\) 的更简单描述。\footnote{编者注：当 \(\cR\) 单连通或伴随时，参见 \SGARefLink{sga3:XXI:7.4.5}{7.4.5}。}

\begin{definition}\label{sga3:XXI:6.7.3}
以 \(\operatorname{Aut}^{s}(\cR)\)\footnote{编者注：上标 \(s\) 用以提示“半单”（semisimple）。}表示满足下图交换的 \(u\in\operatorname{Aut}(\cR)\) 所构成的集合：
\begin{center}
\begin{tikzcd}[column sep=huge,row sep=large]
\cR \arrow[rr,"u"] \arrow[dr]
  && \cR \arrow[dl] \\
& \operatorname{rad}(\cR) &
\end{tikzcd}
\end{center}
记
\[
 E_\Delta^{s}(\cR)
 =E_\Delta(\cR)\cap\operatorname{Aut}^{s}(\cR).
\]
\end{definition}

\begin{remark}\label{sga3:XXI:6.7.4}
若 \(u\in\operatorname{Aut}(\cR)\)，则 \(u\in\operatorname{Aut}^{s}(\cR)\)，当且仅当
\[
 (u-\operatorname{id})(M)\subset\VR(R).
\]
特别地，
\[
 W(\cR)\subset\operatorname{Aut}^{s}(\cR).
\]
\end{remark}

\begin{proposition}\label{sga3:XXI:6.7.5}
对每个单根系 \(\Delta\)，群 \(\operatorname{Aut}^{s}(\cR)\) 为 \(W(\cR)\) 与 \(E_\Delta^{s}(\cR)\) 的半直积，其中前者为正规子群。
\end{proposition}

由函子性，\(\cR\) 的每个自同构诱导 \(\operatorname{ad}(\cR)\) 的一个自同构。因此有典范同态
\[
 \operatorname{Aut}(\cR)\longrightarrow
 \operatorname{Aut}(\operatorname{ad}(\cR)).
\]

\begin{lemma}\label{sga3:XXI:6.7.6}
同态
\[
 \operatorname{Aut}^{s}(\cR)\longrightarrow
 \operatorname{Aut}(\operatorname{ad}(\cR))
\]
为单射。
\end{lemma}

\begin{proof}
设 \(u\) 为 \(M\) 的自同构，使
\[
 (u-\operatorname{id})(M)\subset\VR(R)
\]
且
\[
 {}^tu(\alpha^\ast)=\alpha^\ast
 \qquad(\alpha^\ast\in R^\ast).
\]
对每个 \(x\in M\)，有
\[
 (\alpha^\ast,u(x)-x)
 =({}^tu(\alpha^\ast)-\alpha^\ast,x)=0.
\]
故由 \ref{sga3:XXI:1.2.5}，\(u(x)-x=0\)。
\end{proof}

\begin{lemma}\label{sga3:XXI:6.7.7}
群 \(\operatorname{Aut}^{s}(\cR)\) 有限。
\end{lemma}

\begin{proof}
只需在 \(\cR\) 伴随时证明 \(\operatorname{Aut}(\cR)\) 有限。此时 \(M\) 由 \(R\) 生成，故 \(\cR\) 的每个自同构均由其在 \(R\) 上诱导的置换决定。
\end{proof}
% 底本源行 3759--4107 译毕。

\begin{remark}\label{sga3:XXI:6.7.8}
直接可知，\(\operatorname{Aut}(\cR)\)（分别为 \(E_\Delta(\cR)\)）有限，当且仅当
\[
 \rgred(\cR)-\rgss(\cR)\leq1.
\]
\end{remark}

\subsection{约化根数据的 \(p\)-态射}
\label{sga3:XXI:6.8}

本小节始终固定一个正整数 \(p\)。

\begin{definition}\label{sga3:XXI:6.8.1}
设
\[
 \cR=(M,M^\ast,R,R^\ast),
 \qquad
 \cR'=(M',M'^\ast,R',R'^\ast)
\]
为两个约化根数据。若存在双射
\[
 u:R\xrightarrow{\sim}R'
\]
及映射
\[
 q:R\longrightarrow\{p^n\mid n\in\mathbb N\}
\]
使群同态
\[
 f:M'\longrightarrow M
\]
满足下述条件，则称它为\emph{从 \(\cR'\) 到 \(\cR\) 的 \(p\)-态射}：
\begin{enumerate}
\item[(i)] 对每个 \(\alpha\in R\)，\(f(u(\alpha))=q(\alpha)\alpha\)；
\item[(ii)] 对每个 \(\alpha\in R\)，\({}^tf(\alpha^\ast)=q(\alpha)u(\alpha)^\ast\)。\footnote{编者注：这两个条件蕴含，对一切 \(\alpha,\beta\in R\)，有
  \[
   q(\alpha)\bigl(u(\alpha)^\ast,u(\beta)\bigr)
   =q(\beta)(\alpha^\ast,\beta).
  \]}
\end{enumerate}
\end{definition}

\begin{corollary}\label{sga3:XXI:6.8.2}
\(1\)-态射恰为态射。
\end{corollary}

\begin{corollary}\label{sga3:XXI:6.8.3}
\(p\)-态射的转置为 \(p\)-态射。
\end{corollary}

\begin{lemma}\label{sga3:XXI:6.8.4}
若 \(w\in W(\cR)\) 且 \(\alpha\in R\)，则
\[
 q(w(\alpha))=q(\alpha).
\]
映射
\[
 s_\alpha\longmapsto s_{u(\alpha)}
\]
延拓为同构
\[
 u:W(\cR)\xrightarrow{\sim}W(\cR'),
\]
使
\[
 u(w(\alpha))=u(w)\bigl(u(\alpha)\bigr).
\]
\end{lemma}

\begin{proof}
只需对 \(\alpha,\beta\in R\) 证明
\[
 u(s_\alpha(\beta))=s_{u(\alpha)}u(\beta),
 \qquad
 q(s_\alpha(\beta))=q(\beta).
\]
依次有
\begin{align*}
 f\bigl(s_{u(\alpha)}u(\beta)\bigr)
  &=f(u(\beta))
    -\bigl(u(\alpha)^\ast,u(\beta)\bigr)f(u(\alpha))\\
  &=q(\beta)\beta
    -q(\beta)q(\alpha)^{-1}
       (\alpha^\ast,\beta)q(\alpha)\alpha\\
  &=q(\beta)\bigl(\beta-(\alpha^\ast,\beta)\alpha\bigr)\\
  &=q(\beta)s_\alpha(\beta).
\end{align*}
若
\[
 \gamma=u^{-1}\bigl(s_{u(\alpha)}u(\beta)\bigr),
\]
则
\[
 q(\gamma)\gamma=f(u(\gamma))=q(\beta)s_\alpha(\beta).
\]
根 \(\gamma\) 与 \(s_\alpha(\beta)\) 在 \(\mathbb Q\) 上成比例，故相等或互为相反根。由于 \(q(\gamma)\) 与 \(q(\beta)\) 为正，得到
\[
 q(\gamma)=q(\beta),
 \qquad
 \gamma=s_\alpha(\beta).
\]
\end{proof}

\begin{definition}\label{sga3:XXI:6.8.5}
各数 \(q(\alpha)\) 称为 \(f\) 的\emph{根指数}。
\end{definition}

\begin{example}\label{sga3:XXI:6.8.6}
设 \(\cR\) 为约化根数据，\(q=p^n\)，\(n\in\mathbb N\)。乘以 \(q\) 的映射
\[
 M\longrightarrow M,\qquad x\longmapsto qx,
\]
为 \(p\)-态射，其根指数全等于 \(q\)（且有 \(u=\operatorname{id}\)）。记为
\[
 q:\cR\longrightarrow\cR.
\]
\end{example}

\begin{proposition}\label{sga3:XXI:6.8.7}
采用 \ref{sga3:XXI:6.8.1} 的记号，\(u\) 诱导 \(R\) 的单根系（分别为正根系）的集合到 \(R'\) 的相应集合的同构。
\end{proposition}

\begin{proof}
这直接由 \ref{sga3:XXI:3.1.5}（分别为 \ref{sga3:XXI:3.2.1}）得到。
\end{proof}

\section{结构}\label{sga3:XXI:7}

\subsection{根数据的分解}\label{sga3:XXI:7.1}

\begin{proposition}\label{sga3:XXI:7.1.1}
设 \(\cR\) 为根数据，\(\Delta\) 为单根系。
\begin{enumerate}
\item[\SGAInlineRefTarget{sga3:XXI:7.1.1:item:i}(i)] 设 \(R'\) 与 \(R''\) 为构成 \(R\) 的一个划分的两个闭对称根集。若
  \[
   \Delta'=\Delta\cap R',
   \qquad
   \Delta''=\Delta\cap R'',
  \]
  则
  \[
   R'=R_{\Delta'},
   \qquad
   R''=R_{\Delta''},
  \]
  且 \(\Delta'\) 的每个根都与 \(\Delta''\) 的每个根正交。
\item[\SGAInlineRefTarget{sga3:XXI:7.1.1:item:ii}(ii)] 设 \(\Delta'\) 与 \(\Delta''\) 为构成 \(\Delta\) 的一个划分的两个正交子集。则 \(R_{\Delta'}\) 与 \(R_{\Delta''}\) 构成 \(R\) 的一个划分。
\end{enumerate}
\end{proposition}

\begin{proof}
先证明 \SGARefLink{sga3:XXI:7.1.1:item:i}{(i)}。

\begin{lemma}\label{sga3:XXI:7.1.2}
在 \textup{\SGARefLink{sga3:XXI:7.1.1:item:i}{(i)}} 的假设下，若 \(\alpha\)、\(\beta\) 及 \(\alpha+\beta\) 均为根，则三者全属于 \(R'\)，或全属于 \(R''\)。
\end{lemma}

例如，设 \(\alpha+\beta\in R'\)。不可能有 \(\alpha,\beta\in R''\)，因为 \(R''\) 闭。因此可设 \(\alpha\in R'\)。于是 \(-\alpha\in R'\)，且
\[
 \beta=(\beta+\alpha)-\alpha\in R'.
\]

现对正根 \(\alpha\in R'\cap\mathscr P(\Delta)\) 的阶作归纳，证明 \(R'=R_{\Delta'}\)。若 \(\operatorname{ord}_\Delta(\alpha)=1\)，则
\[
 \alpha\in R'\cap\Delta=\Delta'.
\]
若 \(\operatorname{ord}_\Delta(\alpha)>1\)，则存在 \(\beta\in\Delta\)，使 \(\alpha-\beta\in R\)。由引理，
\[
 \beta\in\Delta',
 \qquad
 \alpha-\beta\in R'.
\]
故由归纳假设，\(\alpha-\beta\in R_{\Delta'}\)，最后有
\[
 \alpha=(\alpha-\beta)+\beta\in R_{\Delta'}.
\]

接着证明 \(\Delta'\) 与 \(\Delta''\) 正交。若 \(\alpha\in\Delta'\) 且 \(\beta\in\Delta''\)，则
\[
 (\beta^\ast,\alpha)\leq0.
\]
若 \((\beta^\ast,\alpha)\neq0\)，则 \(\beta+\alpha\) 为根，与引理矛盾。

现证 \SGARefLink{sga3:XXI:7.1.1:item:ii}{(ii)}。若 \(\Delta'\) 或 \(\Delta''\) 为空，断言直接成立。否则，若 \(R_{\Delta'}\) 与 \(R_{\Delta''}\) 不构成 \(R\) 的划分，则存在形如
\[
 \alpha=\sum_i m_i'\alpha_i'
       +\sum_j m_j''\alpha_j'',
 \qquad
 m_i',m_j''\in\mathbb Z_+,
\]
的根 \(\alpha\)，其中各 \(\alpha_i'\)（分别为 \(\alpha_j''\)）属于 \(\Delta'\)（分别为 \(\Delta''\)）。应用 \ref{sga3:XXI:3.1.2}，必要时交换 \(\Delta'\) 与 \(\Delta''\)，得到关系
\[
 \delta=\gamma+\beta,
 \qquad
 \gamma\in R_{\Delta'},
 \quad
 \beta\in\Delta'',
 \quad
 \delta\in R.
\]
但 \((\beta^\ast,\gamma)=0\)，故由 \ref{sga3:XXI:2.2.5}，\(\gamma-\beta\) 也是根，这不可能。
\end{proof}

\begin{proposition}\label{sga3:XXI:7.1.3}
设 \(\cR\) 为根数据。下述条件等价：
\begin{enumerate}
\item[\SGAInlineRefTarget{sga3:XXI:7.1.3:item:i}(i)] 不存在将 \(R\) 划分为两个闭对称子集的非平凡划分。
\item[\SGAInlineRefTarget{sga3:XXI:7.1:item:ii}(ii)] 对 \(R\) 的某个（分别为每个）单根系 \(\Delta\)，不存在将 \(\Delta\) 划分为两个非空正交子集的划分。
\item[\SGAInlineRefTarget{sga3:XXI:7.1.3:item:iii}(iii)] \(W(\cR)\) 在 \(\VR(R)\) 上的自然表示不可约。
\item[\SGAInlineRefTarget{sga3:XXI:7.1.3:item:iv}(iv)] 对每对根 \((\alpha,\beta)\)，存在根的序列
  \[
   \alpha_0,\alpha_1,\ldots,\alpha_n,
   \qquad
   \alpha_0=\alpha,\quad\alpha_n=\beta,
  \]
  使对 \(i=0,\ldots,n-1\)，\(\alpha_i\) 与 \(\alpha_{i+1}\) 不正交。
\end{enumerate}
\end{proposition}

\begin{proof}
等价关系 \SGARefLink{sga3:XXI:7.1.3:item:i}{(i)}\(\Leftrightarrow\)\SGARefLink{sga3:XXI:7.1:item:ii}{(ii)} 由 \ref{sga3:XXI:7.1.1} 得到，而 \SGARefLink{sga3:XXI:7.1.3:item:iv}{(iv)}\(\Rightarrow\)\SGARefLink{sga3:XXI:7.1.1:item:ii}{(ii)} 直接成立。反之，若 \(\Delta\) 满足 \SGARefLink{sga3:XXI:7.1.1:item:ii}{(ii)}，则当 \(\alpha,\beta\in\Delta\) 时，\SGARefLink{sga3:XXI:7.1.3:item:iv}{(iv)} 成立。每个根与某个单根不正交，例如可由 \ref{sga3:XXI:3.1.1} 得到。此外，\SGARefLink{sga3:XXI:7.1.3:item:iii}{(iii)}\(\Rightarrow\)\SGARefLink{sga3:XXI:7.1.3:item:i}{(i)}：在 \ref{sga3:XXI:7.1.1} 的假设下，空间 \(\VR(R')\) 在 \(W(\cR)\) 下稳定。只需再证明 \SGARefLink{sga3:XXI:7.1.1:item:i}{(i)}\(\Rightarrow\)\SGARefLink{sga3:XXI:7.1.3:item:iii}{(iii)}。

设 \(H\) 为 \(\VR(R)\) 的一个在 \(W(\cR)\) 下稳定的向量子空间。对每个 \(\alpha\in R\)，由等式 \(s_\alpha(H)=H\) 立即得到
\[
 \alpha\in H
 \qquad\text{或}\qquad
 \alpha^\ast\in H^\perp,
\]
其中 \(H^\perp\) 表示 \(H\) 在与 \(\VR(R)\) 对偶的 \(\VR(R^\ast)\) 中的正交补。令
\[
 R'=\{\alpha\in R\mid\alpha\in H\},
 \qquad
 R''=\{\alpha\in R\mid\alpha^\ast\in H^\perp\}.
\]
这给出了 \(R\) 到两个闭对称子集的划分。
\end{proof}

\begin{definition}\label{sga3:XXI:7.1.4}
满足 \ref{sga3:XXI:7.1.3} 的等价条件且半单秩非零的根数据（分别为根系）称为\emph{不可约的}。
\end{definition}

\begin{corollary}\label{sga3:XXI:7.1.5}
对每个根数据 \(\cR\)，存在将 \(R\) 划分为闭对称不可约子集的划分，且在不计次序的意义下唯一。
\end{corollary}

\begin{corollary}\label{sga3:XXI:7.1.6}
每个伴随（分别为单连通）根数据均为不可约伴随（分别为单连通）根数据的积。
\end{corollary}

\begin{proof}
只需考虑伴随情形。此时，断言由下述事实得到：在 \ref{sga3:XXI:7.1.1} 的假设下，
\[
 \Gamma_0(R)=\Gamma_0(R')\times\Gamma_0(R'').
\]
\end{proof}

\begin{corollary}\label{sga3:XXI:7.1.7}
对每个根数据（分别为约化根数据）\(\cR\)，存在同源
\[
 \cR\longrightarrow\cR',
\]
其中 \(\cR'\) 为一个平凡根数据与若干不可约单连通根数据（分别为不可约约化单连通根数据）的积。
\end{corollary}

\subsection{不可约根数据的性质}
\label{sga3:XXI:7.2}

\begin{definition}\label{sga3:XXI:7.2.1}
设 \(\cR\) 为不可约根数据。对每个 \(\alpha\in R\)，令
\[
 \operatorname{long}(\alpha)
 =\ell(\alpha)/\ell(\alpha_0),
\]
其中 \(\alpha_0\in R\) 满足 \(\ell(\alpha_0)\) 最小。数 \(\operatorname{long}(\alpha)\) 称为 \(\alpha\) 的\emph{长度}。
\end{definition}

\begin{lemma}\label{sga3:XXI:7.2.2}
设 \(\cR\) 为不可约根数据。Weyl 群在给定长度的根的集合上传递作用。
\end{lemma}

\begin{proof}
设 \(\alpha,\beta\in R\)。由于 \(W(\cR)\) 在 \(\VR(R)\) 上的表示不可约，\(\alpha\) 不可能与所有 \(w(\beta)\)，\(w\in W(\cR)\)，正交。因此某个 \(w(\beta)\) 与 \(\alpha\) 不正交。由于
\[
 \ell(w(\beta))=\ell(\beta),
\]
断言由 \ref{sga3:XXI:2.3.2} 得到。
\end{proof}

\begin{lemma}\label{sga3:XXI:7.2.3}
若 \(\cR\) 不可约且约化，则
\[
 \operatorname{long}(R)
 \in\bigl\{\{1\},\{1,2\},\{1,3\}\bigr\}.
\]
\end{lemma}

\begin{proof}
由上面使用的观察，对每个 \(\alpha,\beta\in R\)，存在 \(w\in W(\cR)\)，使 \(w(\beta)\) 与 \(\alpha\) 不正交。因此由 \ref{sga3:XXI:2.3.1}，
\[
 \frac{\ell(\alpha)}{\ell(\beta)}
 \in\left\{1,2,3,\frac12,\frac13\right\}.
\]
故 \(\operatorname{long}(\alpha)\) 为 \(1\)、\(2\) 或 \(3\)。但长度 \(2\) 和 \(3\) 不能同时出现，因为其比 \(2/3\) 不在上述集合中。
\end{proof}
% 底本源行 4108--4497 译毕。

\begin{remark}\label{sga3:XXI:7.2.4}
类似的论证给出下述结果。若 \(\cR\) 不可约且非约化，并有 \(\rgss(\cR)>1\)，则
\[
 \operatorname{long}(R)=\{1,2,4\}.
\]
若 \(R_i=\operatorname{long}^{-1}(i)\)，则
\[
 \operatorname{ind}(R)=R_1\cup R_2,\qquad R_4=2R_1,
\]
且 \(R_1\) 中两个不成比例的根互相正交。反之，设 \(R\) 为满足 \(\operatorname{long}(R)=\{1,2\}\) 的不可约约化根系。令 \(R_i=\operatorname{long}^{-1}(i)\)，并假设 \(R_1\) 中两个不成比例的根互相正交。则 \(R\cup2R_1\) 不可约且非约化，且
\[
 \operatorname{ind}(R\cup2R_1)=R.
\]
\end{remark}

\begin{lemma}\label{sga3:XXI:7.2.5}
若 \(\cR\) 为不可约根数据，则 \(\cR^\ast\) 也不可约，且当 \(\alpha\) 遍历 \(R\) 时，乘积
\[
 \operatorname{long}(\alpha)\operatorname{long}(\alpha^\ast)
\]
为常数。
\end{lemma}

\begin{proof}
这直接由 \ref{sga3:XXI:7.1.3}\,\SGARefLink{sga3:XXI:7.2:item:iv}{(iv)} 及 \ref{sga3:XXI:2.2.6} 得到。
\end{proof}

\begin{definition}\label{sga3:XXI:7.2.6}
设 \(\cR\) 为任意根数据。\(\alpha\in R\) 的\emph{长度}，记为 \(\operatorname{long}(\alpha)\)，是 \(\alpha\) 在其不可约分量中的长度。
\end{definition}

\begin{lemma}\label{sga3:XXI:7.2.7}
存在唯一的群同态
\[
 u:\Gamma_0(R)\longrightarrow\Gamma_0(R^\ast)
\]
使
\[
 u(\alpha)=\operatorname{long}(\alpha)\alpha^\ast
 \qquad(\alpha\in R).
\]
\end{lemma}

\begin{proof}
由 \ref{sga3:XXI:3.5.5}，只需验证每当 \(\alpha,\beta,\alpha+\beta\in R\) 时，
\[
 \operatorname{long}(\alpha)\alpha^\ast+
 \operatorname{long}(\beta)\beta^\ast
 =
 \operatorname{long}(\alpha+\beta)(\alpha+\beta)^\ast.
\]
但由 \ref{sga3:XXI:7.1.2}，\(\alpha,\beta,\alpha+\beta\) 属于 \(R\) 的同一不可约分量；故断言化归到 \ref{sga3:XXI:1.2.2}。
\end{proof}

\begin{remark}\label{sga3:XXI:7.2.8}
设 \(u\) 如 \ref{sga3:XXI:7.2.7}。对 \(\alpha,\beta\in R\)，有
\[
 (u(\alpha),\beta)=(u(\beta),\alpha).
\]
事实上，这等价于
\[
 \operatorname{long}(\alpha)(\alpha^\ast,\beta)
 =
 \operatorname{long}(\beta)(\beta^\ast,\alpha),
\]
当 \(\alpha\) 与 \(\beta\) 正交时，这显然成立。若二者不正交，则属于 \(R\) 的同一不可约分量，断言化归到 \ref{sga3:XXI:1.2.1} 的公式 \SGARefLink{sga3:XXI:eq:9}{(9)}。
\end{remark}

\begin{remark}\label{sga3:XXI:7.2.9}
对称双线性型
\[
 (x,y)\longmapsto(u(x),y)
\]
在 \(\Gamma_0(R)\) 上为正且非退化。
\end{remark}

\begin{proof}
设各 \(R_i\) 为 \(R\) 的不可约分量。则
\[
 \Gamma_0(R)=\prod_i\Gamma_0(R_i),
\]
且双线性型 \((u(x),y)\) 为各群 \(\Gamma_0(R_i)\) 上的型
\[
 2\ell(\alpha_i)^{-1}(p(x),y)
\]
之积，其中 \(\ell(\alpha_i)\) 是 \(\alpha\in R_i\) 时 \(\ell(\alpha)\) 的最小值。由 \ref{sga3:XXI:1.2.6}，这些对称双线性型为正且非退化。
\end{proof}

\subsection{Cartan 矩阵}\label{sga3:XXI:7.3}

设 \(\cR\) 为根数据。若 \(\Delta\) 为单根系，则 \(\cR\) 相对于 \(\Delta\) 的 \emph{Cartan 矩阵}是以 \(\Delta\) 为指标集的方阵，定义为
\[
 a_{\alpha,\beta}=(\alpha^\ast,\beta),
 \qquad \alpha,\beta\in\Delta.
\]
若 \(\Delta'\) 为另一个单根系，且 \(w\in W(\cR)\) 满足 \(w(\Delta)=\Delta'\)，则
\[
 (w(\alpha)^\ast,w(\beta))=(\alpha^\ast,\beta).
\]
因此，相对于 \(\Delta'\) 的 Cartan 矩阵由相对于 \(\Delta\) 的矩阵通过 \(w\) 定义的指标集同构 \(\Delta\to\Delta'\) 得到。故在指标集的同构意义下，Cartan 矩阵只依赖于 \(\cR\)。

\begin{proposition}\label{sga3:XXI:7.3.1}
Cartan 矩阵具有下述性质。
\begin{enumerate}
\item[\SGAInlineRefTarget{sga3:XXI:7.3.1:item:i}(i)] \(a_{\alpha,\alpha}=2\)，且当 \(\alpha\neq\beta\) 时，\(a_{\alpha,\beta}\leq0\)。
\item[\SGAInlineRefTarget{sga3:XXI:7.3.1:item:ii}(ii)] \(a_{\alpha,\beta}=0\) 蕴含 \(a_{\beta,\alpha}=0\)。
\item[\SGAInlineRefTarget{sga3:XXI:7.2:item:iii}(iii)] 存在严格正整数 \(m_\alpha=\operatorname{long}(\alpha)\)，使矩阵
  \[
   (m_\alpha a_{\alpha,\beta})
  \]
  对称、为正且非退化。
\item[\SGAInlineRefTarget{sga3:XXI:7.2:item:iv}(iv)] \((a_{\alpha,\beta})_{\alpha,\beta\in\Delta}\) 的主子式，即
  \[
   \det(a_{\alpha,\beta})_{\alpha,\beta\in\Delta'}
   \qquad(\Delta'\subset\Delta),
  \]
  均严格为正。
\item[\SGAInlineRefTarget{sga3:XXI:7.2:item:v}(v)] 有
  \[
   s_\alpha(\beta)=\beta-a_{\alpha,\beta}\alpha,
   \qquad
   s_\alpha(\beta^\ast)
   =\beta^\ast-a_{\beta,\alpha}\alpha^\ast.
  \]
\end{enumerate}
\end{proposition}

\begin{proof}
\SGARefLink{sga3:XXI:7.2:item:v}{(v)} 即为定义；\SGARefLink{sga3:XXI:7.3.1:item:i}{(i)} 由 \ref{sga3:XXI:3.2.11} 得到，\SGARefLink{sga3:XXI:7.3.1:item:ii}{(ii)} 由 \ref{sga3:XXI:2.2.2} 得到，\SGARefLink{sga3:XXI:7.2:item:iii}{(iii)} 由 \ref{sga3:XXI:7.2.9} 得到。\SGARefLink{sga3:XXI:7.2:item:iv}{(iv)} 直接由 \SGARefLink{sga3:XXI:7.2:item:iii}{(iii)} 及
\[
 \det(m_\alpha a_{\alpha,\beta})_{\alpha,\beta\in\Delta'}
 =
 \left(\prod_{\alpha\in\Delta'}m_\alpha\right)
 \det(a_{\alpha,\beta})_{\alpha,\beta\in\Delta'}
\]
得到。
\end{proof}

\begin{proposition}\label{sga3:XXI:7.3.2}
设 \(\cR\) 与 \(\cR'\) 为两个约化单连通（分别为伴随）根数据，\(\Delta\) 与 \(\Delta'\) 为 \(\cR\) 与 \(\cR'\) 的单根系，\(u:\Delta\to\Delta'\) 为同构。以 \((a_{\alpha,\beta})\) 与 \((a'_{\alpha',\beta'})\) 表示相应的 Cartan 矩阵，并假设
\[
 a'_{u(\alpha),u(\beta)}=a_{\alpha,\beta}.
\]
则存在唯一的 \(\cR\) 到 \(\cR'\) 的同构，在 \(\Delta\) 上诱导 \(u\)。
\end{proposition}

\begin{proof}
只需在伴随情形下给出证明。此时 \(M=\Gamma_0(R)\) 与 \(M'=\Gamma_0(R')\) 分别是由 \(\Delta\) 与 \(\Delta'\) 生成的自由 Abel 群。因此存在唯一的群同构 \(M\to M'\)，在 \(\Delta\) 上诱导 \(u\)；仍将它记为 \(u\)。

证明 \(u(R)\subset R'\)。\(R\) 的每个根 \(\alpha\) 均可写为
\[
 \alpha=s_{\alpha_1}\cdots s_{\alpha_n}(\alpha_{n+1}),
 \qquad \alpha_i\in\Delta.
\]
由对 \(u\) 的假设及关系 \ref{sga3:XXI:7.3.1}\,\SGARefLink{sga3:XXI:7.2:item:v}{(v)}，有
\[
 u(\alpha)=
 s_{u(\alpha_1)}\cdots s_{u(\alpha_n)}
 \bigl(u(\alpha_{n+1})\bigr).
\]
只需再证明 \({}^tu(R'^\ast)\subset R^\ast\)。由 \ref{sga3:XXI:1.2.5}，\(M^\ast\)（分别为 \(M'^\ast\)）中的元素由其与 \(R\) 或 \(\Delta\)（分别为 \(R'\) 或 \(\Delta'\)）的对偶配对决定，故得到此结论。
\end{proof}

\begin{corollary}\label{sga3:XXI:7.3.3}
一个约化单连通或伴随根数据在同构意义下由其 Cartan 矩阵决定。
\end{corollary}

\begin{corollary}\label{sga3:XXI:7.3.4}
设 \(\cR\) 为约化单连通（分别为伴随）根数据，\(\Delta\) 为单根系。群 \(E_\Delta(\cR)\) 与保持 Cartan 矩阵的集合 \(\Delta\) 的自同构群相认同。
\end{corollary}

\begin{remark}\label{sga3:XXI:7.3.5}
给定例如满足 \SGARefLink{sga3:XXI:7.3.1:item:i}{(i)}、\SGARefLink{sga3:XXI:7.3.1:item:ii}{(ii)} 及 \SGARefLink{sga3:XXI:7.2:item:iv}{(iv)} 的 Cartan 矩阵，相应根数据的存在性在不使用分类的情况下不易直接解决。
\end{remark}

\subsection{Dynkin 图}\label{sga3:XXI:7.4}

\begin{definition}\label{sga3:XXI:7.4.1}
有限集 \(\Delta\) 上的一个 \emph{Dynkin 图结构}（显然应避免使用“概形”一词），由 \(\Delta\) 中不同元素的某些对所构成的集合，以及一个映射
\[
 \Delta\longrightarrow\{1,2,3\}
\]
组成；这些对称为\emph{相连对}。此种结构之间的同构概念是显然的。
\end{definition}

\begin{definition}\label{sga3:XXI:7.4.2}
设 \(\cR\) 为根数据，\(\Delta\) 为单根系。\(\cR\) 相对于 \(\Delta\) 的 \emph{Dynkin 图}为集合 \(\Delta\)，其中两个单根相连，当且仅当它们不正交，并给每个根赋以其长度。
\end{definition}

\begin{proposition}\label{sga3:XXI:7.4.3}
Dynkin 图与 Cartan 矩阵互相唯一决定。
\end{proposition}

\begin{proof}
等价关系
\[
 \alpha\text{ 与 }\beta\text{ 不相连}
 \quad\Longleftrightarrow\quad
 a_{\alpha,\beta}=0
\]
及关系
\[
 \operatorname{long}(\alpha)a_{\alpha,\beta}
 =
 \operatorname{long}(\beta)a_{\beta,\alpha},
\]
连同图的每个连通分量中 \(\inf\operatorname{long}(\alpha)=1\)，由相连关系及长度决定各 \(a_{\alpha,\beta}\)，反之亦然。细节留待读者验证。
\end{proof}

\begin{corollary}\label{sga3:XXI:7.4.4}
约化单连通或伴随根数据由其 Dynkin 图决定。
\end{corollary}

\begin{corollary}\label{sga3:XXI:7.4.5}
设 \(\cR\) 为约化单连通或伴随根数据，\(\Delta\) 为单根系。群 \(E_\Delta(\cR)\) 与 \(\cR\) 相对于 \(\Delta\) 的 Dynkin 图的自同构群相认同；即与 \(\Delta\) 的保持长度及相连关系的置换群相认同。
\end{corollary}

\begin{remark}\label{sga3:XXI:7.4.6}
通常的方法\footnote{参见 Bourbaki，\emph{Groupes et algèbres de Lie}，第 VI 章，第 4.2 号，或 Sophus Lie Seminar。}
\footnote{编者注：略有不同的证明亦可参见 [De80]。}
对连通 Dynkin 图作出分类，并证明每个图确实对应于一个不可约约化单连通根数据。所得即为熟知的各型：
\begin{center}
\begin{tikzpicture}[
  x=.86cm,y=.78cm,
  dyn/.style={circle,fill=black,inner sep=1.35pt},
  wt/.style={font=\small,above=2pt},
  typ/.style={font=\normalsize,anchor=west}
]
% A_n
\foreach \x in {0,1,2,4,5,6} {
  \node[dyn] (A\x) at (\x,0) {};
  \node[wt] at (A\x) {$1$};
}
\draw (A0)--(A1)--(A2);
\draw (A4)--(A5)--(A6);
\node at (3,0) {$\cdots$};
\node[typ] at (7.3,0) {$A_n,\qquad n\geq1.$};

% B_n
\foreach \x in {0,1,2,4,5} {
  \node[dyn] (B\x) at (\x,-1.55) {};
  \node[wt] at (B\x) {$2$};
}
\node[dyn] (B6) at (6,-1.55) {};
\node[wt] at (B6) {$1$};
\draw (B0)--(B1)--(B2);
\draw (B4)--(B5)--(B6);
\node at (3,-1.55) {$\cdots$};
\node[typ] at (7.3,-1.55) {$B_n,\qquad n\geq2.$};

% C_n
\foreach \x in {0,1,2,4,5} {
  \node[dyn] (C\x) at (\x,-3.10) {};
  \node[wt] at (C\x) {$1$};
}
\node[dyn] (C6) at (6,-3.10) {};
\node[wt] at (C6) {$2$};
\draw (C0)--(C1)--(C2);
\draw (C4)--(C5)--(C6);
\node at (3,-3.10) {$\cdots$};
\node[typ] at (7.3,-3.10) {$C_n,\qquad n\geq3.$};

% D_n
\foreach \x in {0,1,2,4} {
  \node[dyn] (D\x) at (\x,-4.65) {};
  \node[wt] at (D\x) {$1$};
}
\node[dyn] (Du) at (5,-4.25) {};
\node[wt] at (Du) {$1$};
\node[dyn] (Dl) at (5,-5.05) {};
\node[wt] at (Dl) {$1$};
\draw (D0)--(D1)--(D2);
\draw (D4)--(Du);
\draw (D4)--(Dl);
\node at (3,-4.65) {$\cdots$};
\node[typ] at (7.3,-4.65) {$D_n,\qquad n\geq4.$};

% E_n
\foreach \x in {0,1,3,4,5} {
  \node[dyn] (E\x) at (\x,-6.35) {};
  \node[wt] at (E\x) {$1$};
}
\node[dyn] (Ev) at (3,-7.15) {};
\node[right=2pt,font=\small] at (Ev) {$1$};
\draw (E0)--(E1);
\draw (E3)--(E4)--(E5);
\draw (E3)--(Ev);
\node at (2,-6.35) {$\cdots$};
\node[typ] at (7.3,-6.65) {$E_n,\qquad n=6,7,8.$};

% F_4
\foreach \x/\w in {0/1,1/1,2/2,3/2} {
  \node[dyn] (F\x) at (\x,-8.15) {};
  \node[wt] at (F\x) {$\w$};
}
\draw (F0)--(F1)--(F2)--(F3);
\node[typ] at (7.3,-8.15) {$F_4.$};

% G_2
\node[dyn] (G0) at (2,-9.70) {};
\node[wt] at (G0) {$1$};
\node[dyn] (G1) at (3,-9.70) {};
\node[wt] at (G1) {$3$};
\draw (G0)--(G1);
\node[typ] at (7.3,-9.70) {$G_2.$};
\end{tikzpicture}
\end{center}
由 \ref{sga3:XXI:7.4.5}，立即可得相应的群 \(E_\Delta(\cR)\)：
\[
\begin{array}{ll}
E_\Delta(\cR)=\{e\},
 & A_1,\ B_n,\ C_n,\ E_7,\ E_8,\ F_4,\ G_2,\\[2pt]
E_\Delta(\cR)=\mathbb Z/2\mathbb Z,
 & A_n\ (n\geq2),\ D_n\ (n\geq5),\ E_6,\\[2pt]
E_\Delta(\cR)=\mathfrak S_3,
 & D_4.
\end{array}
\]
\end{remark}
% 底本源行 4498--4908 译毕。

\subsection{关于 \(p\)-态射的补充}\label{sga3:XXI:7.5}

设 \(f:\cR\to\cR'\) 为 \(p\)-态射（参见 \ref{sga3:XXI:6.8}）。由定义显然可见，与 \(f\) 相伴的双射 \(u:R\xrightarrow{\sim}R'\) 使 \(R\) 与 \(R'\) 的单根系、正根系、不可约分量等等互相对应。因此为简单起见，可设 \(R\) 与 \(R'\) 不可约。

\begin{lemma}\label{sga3:XXI:7.5.1}
若 \(R\) 与 \(R'\) 不可约，则存在 \(k\in\mathbb Q\)，使对每个 \(\alpha\in R\)，有
\[
 k\,\operatorname{long}(u(\alpha))
 =
 q(\alpha)^2\operatorname{long}(\alpha).
\]
\end{lemma}

\begin{proof}
有
\[
 \operatorname{long}(\alpha)(\alpha^\ast,\beta)
 =
 \operatorname{long}(\beta)(\beta^\ast,\alpha),
\]
且类似地，有
\[
 \operatorname{long}(u(\alpha))
   (u(\alpha)^\ast,u(\beta))
 =
 \operatorname{long}(u(\beta))
   (u(\beta)^\ast,u(\alpha)).
\]
立即推出，对不正交的 \(\alpha,\beta\)，有
\[
 \frac{q(\alpha)^2\operatorname{long}(\alpha)}
      {\operatorname{long}(u(\alpha))}
 =
 \frac{q(\beta)^2\operatorname{long}(\beta)}
      {\operatorname{long}(u(\beta))}.
\]
结论遂由 \ref{sga3:XXI:7.1.3}\,\SGARefLink{sga3:XXI:7.1.3:item:iv}{(iv)} 得到。
\end{proof}

\begin{remark}\label{sga3:XXI:7.5.2}
由 \ref{sga3:XXI:7.2.2} 及 \ref{sga3:XXI:6.8.4}，\(q(\alpha)\) 只依赖于 \(\operatorname{long}(\alpha)\)。于是容易看出，若 \(q(\alpha)\) 非常值，则乘积 \(q(\alpha)\operatorname{long}(\alpha)\) 为常值；故 \(p=2\) 或 \(3\)。考察上述各图，得到四种可能的情形（同一字母表示 Dynkin 图及其相应的约化单连通根数据）：
\[
\begin{array}{llll}
p=2, & B_n\xrightarrow{\,f_1\,}C_n,
     & C_n\xrightarrow{\,f_2\,}B_n,
     & (C_2=B_2),\\
p=2, & F_4\xrightarrow{\,g\,}F_4,\\
p=3, & G_2\xrightarrow{\,h\,}G_2.
\end{array}
\]
读者可注意到，\(f_1\circ f_2\)、\(f_2\circ f_1\)、\(g\circ g\) 及 \(h\circ h\) 均为 \ref{sga3:XXI:6.8.6} 所述形式的 \(p\)-态射。
\end{remark}

\medskip
\noindent\textbf{7.5.3。}\label{sga3:XXI:7.5.3}
上述描述立即表明：若 \(\cR\) 与 \(\cR'\) 为两个半单秩至多为 \(2\) 的约化根数据，且存在从 \(\cR'\) 到 \(\cR\) 的 \(p\)-态射，则 \(\cR\) 与 \(\cR'\) 的型相同。更确切地，有下表。

\medskip
\noindent\emph{记号。}
设 \(f:\cR'\to\cR\) 为 \(p\)-态射。符号 \(q\) 与 \(q_1\) 表示 \(p\) 的任意正幂。对秩二根系，采用第 \ref{sga3:XXI:4} 节的记号，其中 \(\alpha,\beta\) 为单根，且 \(\ell(\alpha)\leq\ell(\beta)\)。

\begin{center}
\begingroup
\renewcommand{\arraystretch}{1.35}
\resizebox{\linewidth}{!}{%
\begin{tabular}{|c|c|c|c|}
\hline
\textbf{型} & \(p\) & \textbf{\(f\) 的值}
 & \textbf{\({}^tf\) 的值}\\
\hline
平凡 & 任意 & --- & ---\\
\hline
\(A_1\) & 任意
 & \(f(\alpha')=q\alpha\)
 & \({}^tf(\alpha^\ast)=q\alpha'^\ast\)\\
\hline
\(A_1\times A_1\) & 任意
 & \(\begin{array}{c}f(\alpha')=q\alpha\\
                      f(\beta')=q_1\beta\end{array}\)
 & \(\begin{array}{c}{}^tf(\alpha^\ast)=q\alpha'^\ast\\
                      {}^tf(\beta^\ast)=q_1\beta'^\ast\end{array}\)\\
\hline
\(A_2,B_2,G_2\) & 任意
 & \(\begin{array}{c}f(\alpha')=q\alpha\\
                      f(\beta')=q\beta\end{array}\)
 & \(\begin{array}{c}{}^tf(\alpha^\ast)=q\alpha'^\ast\\
                      {}^tf(\beta^\ast)=q\beta'^\ast\end{array}\)\\
\hline
\(B_2\) & \(p=2\)
 & \(\begin{array}{c}f(\alpha')=q\beta\\
                      f(\beta')=2q\alpha\end{array}\)
 & \(\begin{array}{c}{}^tf(\alpha^\ast)=2q\beta'^\ast\\
                      {}^tf(\beta^\ast)=q\alpha'^\ast\end{array}\)\\
\hline
\(G_2\) & \(p=3\)
 & \(\begin{array}{c}f(\alpha')=q\beta\\
                      f(\beta')=3q\alpha\end{array}\)
 & \(\begin{array}{c}{}^tf(\alpha^\ast)=3q\beta'^\ast\\
                      {}^tf(\beta^\ast)=q\alpha'^\ast\end{array}\)\\
\hline
\end{tabular}}
\endgroup
\end{center}

\section*{参考文献}
\addcontentsline{toc}{section}{参考文献}
\label{sga3:XXI:bibliography}

\begin{description}
\item[{[BLie]}] N.~Bourbaki，
  \emph{Groupes et algèbres de Lie}，第 IV--VI 章，
  Hermann，1968。
\item[{[De80]}] M.~Demazure，
  ``A, B, C, D, E, F, etc.,'' 见
  \emph{Séminaire sur les singularités des surfaces}
  （Palaiseau，1976--1977），第 221--227 页，M.~Demazure、
  H.~C.~Pinkham 及 B.~Teissier 编，
  Lecture Notes in Mathematics \textbf{777}，
  Springer-Verlag，1980。
\end{description}
% 底本源行 4909--5057 译毕；本分工范围首译完成。
